\begin{document}
\begin{center}
{\Large\textbf{Масштаб, жидкость и изменение поверхности}}\\[2mm]
{\small ЕГЭ, стереометрия, задание №3 \quad | \quad Фамилия\ \answerline{6cm}}
\end{center}
\ifteacher
\begin{center}\fbox{\textbf{Учительская версия:} ответы и краткие опоры выделены курсивом.}\end{center}
\fi

\textbf{Цель.} До подстановки чисел назвать, какая степень коэффициента нужна:
первая для длины, вторая для площади, третья для объёма.

\section*{1. Изменение размеров цилиндра}
\stereotask{1}{Объём первого цилиндра равен \(5\). У второго цилиндра высота в
\(2{,}5\) раза меньше, а радиус основания в \(3\) раза больше. Найдите объём
второго цилиндра.}{\csname figStereoV19\endcsname}
{Коэффициент объёма \(3^2/2{,}5=3{,}6\). Ответ: \(18\).}

\stereotask{2}{В цилиндрическом сосуде уровень жидкости достигает \(25\) см.
На какой высоте будет находиться уровень жидкости во втором сосуде, если его
диаметр в \(2{,}5\) раза больше диаметра первого?}{\csname figStereoV20\endcsname}
{Площадь основания увеличилась в \(2{,}5^2\) раза, поэтому высота равна \(25/2{,}5^2=4\).}

\section*{2. Подобные конусы и вытеснение}
\stereotask{3}{В сосуде в форме конуса уровень жидкости достигает \(0{,}25\)
высоты. Объём жидкости равен \(5\) мл. Сколько миллилитров жидкости нужно
долить, чтобы полностью наполнить сосуд?}{\csname figStereoV17\endcsname}
{Жидкость составляет \(0{,}25^3=1/64\) объёма сосуда. Полный объём \(320\) мл,
долить нужно \(315\) мл.}

\stereotask{4}{Площадь боковой поверхности конуса равна \(30\). Сечение,
параллельное основанию, делит высоту в отношении \(2:3\), считая от вершины
конуса. Найдите площадь боковой поверхности отсечённого конуса.}
{\csname figStereoV35\endcsname}
{Линейный коэффициент малого конуса \(2/5\), площадь уменьшается в \(25/4\) раза.
Ответ: \(30\cdot(2/5)^2=4{,}8\).}

\stereotask{5}{В сосуд правильной треугольной призмы налили
\(1100\ \text{см}^3\) воды и полностью погрузили деталь. Уровень жидкости
поднялся с \(25\) см до \(29\) см. Найдите объём детали.}
{\csname figStereoV18\endcsname}
{Площадь основания сосуда \(1100/25=44\), прирост уровня \(4\). Объём детали
\(44\cdot4=176\).}

\section*{3. Новая поверхность и сумма площадей}
\stereotask{6}{Из единичного куба вырезана правильная четырёхугольная призма
со стороной основания \(0{,}4\) и боковым ребром \(1\). Найдите площадь
поверхности оставшейся части куба.}{\csname figStereoV09\endcsname}
{Убраны два квадрата \(0{,}4^2\), добавлена боковая поверхность \(4\cdot0{,}4\).
Ответ: \(6-2\cdot0{,}16+1{,}6=7{,}28\).}

\stereotask{7}{Из единичного куба вырезана правильная четырёхугольная призма
со стороной основания \(0{,}6\) и боковым ребром \(1\). Найдите площадь
поверхности оставшейся части куба.}{\csname figStereoV10\endcsname}
{Убраны два квадрата \(0{,}6^2\), добавлена боковая поверхность \(4\cdot0{,}6\).
Ответ: \(6-2\cdot0{,}36+2{,}4=7{,}68\).}

\stereotask{8}{Радиусы двух шаров равны \(7\) и \(24\). Найдите радиус шара,
площадь поверхности которого равна сумме площадей поверхностей двух данных
шаров.}{\csname figStereoV13\endcsname}
{Складываются квадраты радиусов: \(R=\sqrt{7^2+24^2}=25\).}

\ifteacher
\section*{Короткая проверка урока}
Ученик готов к смешанной работе, если сначала называет степень коэффициента и
не путает площадь основания сосуда с объёмом воды.
\fi
\end{document}

